\chapter*{คำนำ}
\addcontentsline{toc}{chapter}{คำนำ}

คู่มือวิชาการและชุดพรอมพต์การเรียนรู้เล่มนี้ จัดทำขึ้นเพื่อเป็นแนวทางเชิงปฏิบัติการสำหรับนักศึกษา นักวิจัย เกษตรกรรุ่นใหม่ และผู้สนใจเริ่มต้นศึกษาการพัฒนานวัตกรรมเกษตรดิจิทัล โดยมุ่งเน้นการตรวจวัดและวิเคราะห์ค่าความเป็นกรด-ด่างของดิน (Soil pH) ซึ่งเป็นตัวแปรควบคุมหลักที่มีความสำคัญยิ่งยวดต่อความสามารถในการละลายและการดูดซึมธาตุอาหารของพืชเศรษฐกิจ เช่น ทุเรียนหมอนทองในแถบภาคตะวันออก

ในการตรวจวัดค่าความเป็นกรด-ด่างของดินในสภาพแปลงเกษตรกรรมจริง สัญญาณไฟฟ้าเคมีที่ได้จากหัววัดมักประสบปัญหาความคลาดเคลื่อนอย่างรุนแรงจากความผันผวนของสภาพแวดล้อม โดยเฉพาะการเลื่อนลอยตามอุณหภูมิดิน (Soil Temperature Drift) ซึ่งส่งผลกระทบต่อความชันศักย์ไฟฟ้าตามสมการเนิร์นสต์ (Nernstian Slope) และการแตกตัวของน้ำ ($pK_w$) ตลอดจนผลกระทบจากปริมาณความชื้นในดิน (Soil Moisture Effect) ซึ่งก่อให้เกิดการเปลี่ยนแปลงของอิมพีแดนซ์รอยต่อของเหลว (Liquid Junction Impedance) ในสภาวะดินแห้ง

เพื่อแก้ไขปัญหาเชิงโครงสร้างดังกล่าว คู่มือเล่มนี้ได้นำเสนอการพัฒนาแอปพลิเคชันสมาร์ตโฟน \textbf{SoilpH-HT} บนระบบปฏิบัติการแอนดรอยด์ โดยผสานขุมพลังของโครงข่ายประสาทเทียมอิงฟิสิกส์ (\textbf{Physics-Informed Neural Network} หรือ \textbf{PINN}) ซึ่งทำการฝังทฤษฎีเคมีไฟฟ้าและอุณหพลศาสตร์ลงในโครงข่ายประสาทเทียม เพื่อทำการชดเชยความคลาดเคลื่อนร่วมระหว่างอุณหภูมิและความชื้นแบบไม่เป็นเชิงเส้นในระดับอุปกรณ์ (Edge AI) ได้อย่างแม่นยำภายในเวลาต่ำกว่า 0.2 มิลลิวินาที พร้อมเชื่อมต่อกับหัววัด Soil Parameter Tester ผ่านพอร์ต USB Type-C OTG โดยตรง

นอกจากนี้ คู่มือยังครอบคลุมการออกแบบส่วนต่อประสานผู้ใช้ที่ปรับขนาดอัตโนมัติตามขนาดหน้าจอสมาร์ตโฟน (Responsive UI Auto-Scaling) การพัฒนาระบบกล้องปัญญาประดิษฐ์เพื่อบันทึกภาพถ่ายดินพร้อมพิกัดดาวเทียมและข้อมูลเทเลเมทรีบนภาพ (Telemetry HUD Overlay) และการแปลผลปฐพีวิทยาเพื่อออกใบสั่งปูนโดโลไมต์ปรับปรุงโครงสร้างดินตามลักษณะเนื้อดินจริง

คณะผู้จัดทำหวังเป็นอย่างยิ่งว่า คู่มือวิชาการและชุดพรอมพต์วิศวกรรมเล่มนี้ จะเป็นประโยชน์ต่อการจัดการเรียนการสอน การวิจัยเชิงลึก และการพัฒนาเทคโนโลยีเกษตรแม่นยำของประเทศให้ก้าวสู่ระดับสากลอย่างยั่งยืนสืบไป

\vspace{1.5cm}
\begin{flushright}
    ผู้ช่วยศาสตราจารย์ ดร.ชีวะ ทัศนา\\
    สาขาวิชาฟิสิกส์ คณะวิทยาศาสตร์และเทคโนโลยี\\
    มหาวิทยาลัยราชภัฏรำไพพรรณี
\end{flushright}
